\def\topic{Differential geometry}
\ifdefined\MyHeaderIncluded
\else
\def\qwertyflag{true}

\ifdefined\MyHeaderIncluded

\else
  \newcommand{\foobar}{FooBar}
  \def\MyHeaderIncluded{true}
% if is till the end of the file

\documentclass[10pt,a4paper]{article}
\usepackage[utf8]{inputenc}
%\usepackage[german]{babel}
%\usepackage[cyrillic]{babel}
\usepackage[T1]{fontenc}
%\usepackage[T2A]{fontenc}
%\usepackage[english, russian]{babel}
\usepackage{amsmath}
\usepackage{amsfonts}
\usepackage{amssymb}
\usepackage{amsthm}
\usepackage{mathtools} %coloneqq

\usepackage{bussproofs} %

%\renewcommand{\thesubsection}{\arabic{subsection}}
%\newtheorem{thm}{Theorem}[section]
%\renewcommand{\thethm}{\arabic{subsection}.\arabic{thm}}

\usepackage{thmtools}
\usepackage{hyperref}
\hypersetup{
    colorlinks=true,
    linkcolor=blue,
    filecolor=magenta,      
    urlcolor=cyan,
    pdftitle={Overleaf Example},
    %pdfpagemode=FullScreen,
    }
%\usepackage{stix-type1}%\turnediota
\theoremstyle{definition}

\newtheorem{constructor}{Constructor}[section]
\newtheorem{iaxiom}{Axiom}
\newtheorem{axiom}{Axiom}
\newtheorem*{axiom*}{Axiom}
\newtheorem{definition}{Definition}[section]
\newtheorem*{definition*}{Definition}
\newtheorem{remark}{Remark}[section]
\newtheorem*{remark*}{Remark}
\newtheorem{recall}{Recall}[section]
\newtheorem*{recall*}{Recall}

\newtheorem{lemma}{Lemma}[section]
\newtheorem{theorem}{Theorem}[section]
\newtheorem*{theorem*}{Theorem}
\newtheorem{note}{Note}[section]
\usepackage{etoolbox}

%\newcommand*{\@Dashv}[2]{%
%  \reflectbox{$\m@th#1#2$}%
%}
\newcommand*{\separation}[3]{% 1)binder, 2)set, 3)condition
  \{#1\in #2: #3\}
}

% https://www.freecodecamp.org/news/simple-regex-tricks-for-beginners-3acb3fa257cb/
%  \\mathcal\{T\} |---> \topology
%  \\mathcal\{([^\}]*)\}   |--->  $1
\newcommand{\AND}{\ \ \&\ \ }
\newcommand{\any}{\mathrm{\ any}}
\newcommand{\borel}{\mathcal{B}}
\newcommand{\can}{\mathrm{can}}
\newcommand{\canonize}{\mathrm{canonize}}
\newcommand{\closure}{\mbox{cl}}
\newcommand{\complex}{\mathbb{C}}
\newcommand{\converges}{\longrightarrow}
\newcommand{\countable}{\mbox{\ countable}}
\newcommand{\ldefined}{\mbox{\ defined}}
\newcommand{\dom}{\mbox{dom}}
\newcommand{\Dom}{\mbox{Dom}}
\newcommand{\Done}{\noindent \textbf{Done.}\\}
\newcommand{\lexists}{\mbox{\ exists}}
\newcommand{\lforall}{\mbox{\ forall}}
\newcommand{\lunique}{\mbox{\ unique}}
\newcommand{\family}{\mathrm{\ family}}
\newcommand{\Fm}{\mathrm{Fm}}
\newcommand{\FuncSymb}{\mathrm{FuncSymb}}
\newcommand{\function}{\mathrm{\ function}}
\newcommand{\functional}{\mbox{\ functional}}
\newcommand{\IFF}{\ \Longleftrightarrow\ }
\renewcommand{\iff}{\ \longleftrightarrow\ }
\newcommand{\IMPLIES}{\ \Longrightarrow\ }
\renewcommand{\implies}{\ \longrightarrow\ }
\newcommand{\Inf}{\mathrm{Inf}}
\newcommand{\injective}{\mbox{\ injective}}
\newcommand{\injection}{\mbox{\ injection}}
\newcommand{\interior}{\mbox{int}}
\newcommand{\intersection}{\bigcap}
\newcommand{\isempty}{\mathrm{\ empty}}
\newcommand{\LB}{\mathrm{LB}}
\newcommand{\nghbset}{\mathcal{N}}
\newcommand{\nonempty}{\mathrm{\ nonempty}}
\newcommand{\Ord}{\mathrm{Ord}}
\newcommand{\power}{\mathcal{P}}
\newcommand{\pr}{\mathrm{pr}}
\newcommand{\fst}{\mathrm{fst}}
\newcommand{\snd}{\mathrm{snd}}
\newcommand{\PredSymb}{\mathrm{PredSymb}}
\newcommand{\Proof}{\noindent \textbf{Proof.}\\}
\newcommand{\ran}{\mathrm{ran}}
\newcommand{\Ran}{\mathrm{Ran}}
\newcommand{\real}{\mathbb{R}}
\newcommand{\relation}{\mathrm{\ relation}}
\newcommand{\set}{\mathrm{\ set}}
\newcommand{\singleton}{\mathrm{\ singleton}}
\newcommand{\semantics}{\mathrm{semantics}}
\newcommand{\semialgebra}{\mathcal{S}}
\newcommand{\Sup}{\mathrm{Sup}}
\newcommand{\surjective}{\mbox{\ surjective\ on\ }}
\newcommand{\Term}{\mathrm{Term}}
\newcommand{\Type}{\mathrm{Type}}
\newcommand{\topology}{\mathcal{T}}
\newcommand{\total}{\mbox{\ total\ on\ }}
\newcommand{\transitive}{\mathrm{\ transitive}}
\newcommand{\union}{\bigcup}
\newcommand{\lTrue}{\mathrm{True}}
\newcommand{\lFalse}{\mathrm{False}}
\newcommand{\lUncertain}{\mathrm{Uncertain}}
\newcommand{\lif}{\mathrm{if}}
\newcommand{\lthen}{\mathrm{then}}
\newcommand{\lelse}{\mathrm{else}}
\newcommand{\lexception}{\mathrm{exception}}
\newcommand{\lundef}{\mathrm{undef}}
\newcommand{\univocal}{\mathrm{\ univocal}}
\newcommand{\univocalat}{\mathrm{\ univocal\ at\ }}
\newcommand{\UB}{\mathrm{UB}}
\newcommand{\Un}{\mathcal{U}}
\newcommand{\Var}{\mathrm{Var}}

\newcommand{\notvarphi}{\phi}
\renewcommand{\phi}{\varphi}

\newcommand{\comment}[1]{\quad\mbox{(#1)}}
% local commands
%\preto\section{\undef\blipo} %etoolbox
\preto\subsection{\undef\blipo} %etoolbox
% https://tex.stackexchange.com/questions/105684/limit-macro-newcommand-scope-to-section
\usepackage{graphicx}
\usepackage{authblk}

\setlength\parindent{0pt} % whole doc of \noindent
\usepackage{soul}
%\usepackage{mathgl}

\fi % guard
\ifdefined\topic
\else
\def\topic{No Titel}
\fi
\title{\topic}
\author[1]{Georgy Dunaev}
\begin{document}
\maketitle
\fi
\section{Local Euclidean structure and connectedness}
Here a manifold is a topological space locally homeomorphic to open
subsets of $\mathbb R^n$, for a fixed finite $n\geq0$. The results below
use only this local Euclidean property: neither differentiability nor
the usual additional Hausdorff and second-countability axioms enter
their proofs. Connectedness is a property of the manifold's topology
and does not depend on the choice of charts.

\begin{lemma}[Small path-connected coordinate neighbourhoods]
Every point of a manifold has a base of open path-connected
neighbourhoods.
\end{lemma}
\begin{proof}
Let $p\in M$ and let $O$ be an open neighbourhood of $p$. Choose a
chart $\varphi:U\to V\subseteq\mathbb R^n$ with $p\in U$. The set
$\varphi[U\cap O]$ is open and contains $\varphi(p)$, so it contains
an open Euclidean ball $B$ centred at $\varphi(p)$. The open set
$\varphi^{-1}[B]\subseteq O$ is path connected: the straight segment
between any two points of $B$ stays in $B$, and its composition with
$\varphi^{-1}$ is a path in the manifold. This also works for $n=0$,
when $B$ is the one-point space.
\end{proof}

\begin{theorem}[Components of a manifold]
The connected components of a manifold are open and are exactly its
path components. In particular, a nonempty connected manifold is
path connected.
\end{theorem}
\begin{proof}
Paths define an equivalence relation: constant paths give reflexivity,
reversing the parameter gives symmetry, and concatenating two paths
on the two halves of $[0,1]$ gives transitivity. Fix a path component
$C$. For each $q\in C$, a path-connected open neighbourhood of $q$
lies in $C$, since its paths concatenate with a path from a fixed
point of $C$ to $q$. Thus every path component is open. Their partition
also makes every path component closed, as its complement is the
union of the others.

A path-connected set is connected: a separation would pull back
along a path joining its two sides to a separation of $[0,1]$.
Conversely, any connected subset meeting $C$ must lie in $C$, since
the relatively open sets obtained from $C$ and its complement would
otherwise separate it. Hence $C$ is precisely the maximal connected
subset containing any of its points. These are the connected
components, which proves both assertions.
\end{proof}

\begin{corollary}[Locally constant data on a connected manifold]
If $M$ is nonempty and connected and $f:M\to D$ is locally constant,
then $f$ is constant. Equivalently, every continuous map from $M$ to
a discrete space has one-point image.
\end{corollary}
\begin{proof}
Every fibre of a locally constant map is open, and its complement is
the union of the other open fibres. Choose $p\in M$. The nonempty
fibre of $f(p)$ is therefore both open and closed, so connectedness
makes it all of $M$. A map into a discrete space is continuous exactly
when all its fibres are open, which is exactly local constancy. In
fact this argument applies to every connected topological space.
\end{proof}

\begin{remark}[Connected zero-dimensional manifolds]
Every nonempty connected zero-dimensional manifold $M$ consists of a
single point. Indeed, a coordinate neighbourhood is homeomorphic to the
one-point space $\mathbb{R}^0$, so every singleton in $M$ is open. Thus $M$ is
discrete, and a discrete space with at least two points is disconnected.
The nonemptiness hypothesis excludes the empty manifold if it is allowed
by the chosen convention.
\end{remark}
\ifdefined\qwertyflag
\undef\qwertyflag
\end{document}
\fi
